% Generated draft: requires mathematical and source review.
\documentclass[11pt]{article}
\usepackage[margin=25mm]{geometry}
\usepackage{fontspec}
\setmainfont{Latin Modern Roman}
\usepackage{amsmath,amssymb,mathtools,unicode-math}
\setmathfont{Latin Modern Math}
\usepackage{enumitem,needspace,xurl,hyperref}
\setlength{\parindent}{0pt}
\setlength{\parskip}{5pt}
\setlength{\emergencystretch}{3em}
\newcommand{\N}{\mathbb{N}}
\newcommand{\Z}{\mathbb{Z}}
\newcommand{\Q}{\mathbb{Q}}
\newcommand{\R}{\mathbb{R}}
\newcommand{\C}{\mathbb{C}}
\newcommand{\F}{\mathbb{F}}
\newcommand{\A}{\mathbb{A}}
\newcommand{\PP}{\mathbb{P}}
\newcommand{\E}{\mathbb{E}}
\newcommand{\eps}{\varepsilon}
\newcommand{\dd}{\,\mathrm{d}}
\begin{document}
\section*{UnsolvedMath ID: 30006401}
\textbf{Source problem code:} catalogue entry
\subsection{Definitions and mathematical statement}
For an integer \(k\geq 3\), let \(R(3,k)\) be the least integer \(N\) such that every finite simple graph on \(N\) vertices contains a triangle or an independent set of size \(k\). Write \(\alpha(G)\) for the independence number of a graph \(G\), and use the convention that \(\log\) denotes the natural logarithm.

For a positive integer \(d\), let \(V_d=\F_2^d\). A set \(S\subseteq V_d\) is sum-free if there are no \(a,b,c\in S\), with repetitions allowed, satisfying \(a+b=c\). In particular, \(0\notin S\). Define the Cayley-sum graph \(G_d(S)\) on vertex set \(V_d\) by joining distinct vertices \(x,y\) exactly when \(x+y\in S\). If \(S\) is sum-free, then \(G_d(S)\) is triangle-free. Thus \(\alpha(G_d(S))<k\) implies \(R(3,k)>2^d\).

\textbf{Construction question, conditional on specifying the process.} Can the sum-free process intended by the source produce sum-free sets \(S_j\subseteq\F_2^{d_j}\), for integers \(d_j\to\infty\) and \(k_j\to\infty\), such that
\[
\alpha\bigl(G_{d_j}(S_j)\bigr)<k_j,
\qquad
2^{d_j}\geq\left(\frac12-\varepsilon_j\right)\frac{k_j^2}{\log k_j},
\qquad
\varepsilon_j\geq0,\quad \varepsilon_j\longrightarrow0?
\]
These conditions express an existential construction yielding the requested Ramsey lower bound along an infinite sequence. They do not specify a probability of success. The supplied record does not define the sum-free process, its selection probabilities, its stopping rule, or whether the intended assertion is existence, positive probability, or probability tending to one. Consequently, this is a conditional formulation of the target, not a complete specification of the proposed random construction.

\textbf{Scope and source qualifications.} The original excerpt also appears to state the broader conjecture
\[
R(3,k)=\left(\frac12+o(1)\right)\frac{k^2}{\log k}
\qquad (k\to\infty),
\]
but its displayed formula is corrupted in the supplied transcription. The clean statement asks only about the proposed Cayley-sum construction and a lower bound along an infinite sequence; it does not request a matching upper bound.

The excerpt separately describes sampling a random graph on \(n/(\log n)^2\) vertices at density \(p_0=\alpha_0(\log n/n)^{1/2}\), removing triangles, taking a blow-up, and running a modified triangle-free process. It describes manually steering that process, attributes a simpler proof of an earlier lower bound to the version without the seed step, and suggests that combining randomness with structure may improve on the natural triangle-free process. These are supplied contextual claims. The excerpt supplies neither the parameter restrictions and rounding conventions nor the cleaning, blow-up, steering, and stopping rules needed to define that construction, and it does not explain its relationship to the proposed sum-free process. It expressly says that the authors were unable to provide a construction attaining the proposed constant \(1/2\).
\subsection{Short English statement}
The question asks whether a specified sum-free process on \(\F_2^d\) can produce triangle-free Cayley-sum graphs with sufficiently small independent sets to witness the Ramsey lower bound with asymptotic constant \(1/2\) along infinitely many values of \(k\). The process and its intended success probability are missing. The supplied open status and dated literature assessment require independent review.
\subsection{Sources}
\begin{itemize}
\item Local UnsolvedMath registry, \texttt{data/unsolvedmath/problems.json}, permanent ID \(30006401\), source problem code \texttt{catalogue entry}, title \emph{Cayley-Sum Constructions for Triangle Ramsey Numbers}. The supplied record lists proposed year 2025 and no named proposer.
\item Supplied source citation: \emph{Combinatorics, Probability and Computing (2026)}, Oberwolfach Reports, DOI 10.4171/owr/2025/42, pp. 2243--2322. \url{https://doi.org/10.4171/owr/2025/42}. The literature-source list labels this same reference \emph{Oberwolfach Report 42/2025}.
\item Supplied literature source: \emph{Improving R(3,k) in just two bites}. \url{https://arxiv.org/abs/2510.19718}.
\item Literature assessment dated exactly \texttt{2026-08-22}: the supplied curation labels the problem open and states that the conjectured asymptotic for \(R(3,k)\) remains unresolved, with a lower-bound constant of one half in the 2025--2026 literature but no matching upper bound. This is supplied dated curation, not independently verified. It does not establish that the specific Cayley-sum construction achieves that constant.
\item The original excerpt cites references [3] and [8], Theorem 1, and a lower bound numbered (1), without supplying their identities or complete statements. The registry labels the statement \texttt{corrected\_verified} and describes its changes as typographical; this supplied label does not constitute independent mathematical verification.
\end{itemize}
\end{document}
